\onehalfspacing

\renewcommand{\contentsname}{\centerline{\large СОДЕРЖАНИЕ}}
\tableofcontents

\newpage
\begin{center}
  \textbf{\large ВВЕДЕНИЕ}
\end{center}
\addcontentsline{toc}{chapter}{ВВЕДЕНИЕ}


\textbf{Актуальность темы исследования.} Цифровая трансформация сферы услуг меняет способы взаимодействия предприятий общественного питания с гостями: контакт с заведением начинается до визита, продолжается во время обслуживания и сохраняется после него. Исследования ресторанного бизнеса и малых предприятий показывают рост значения цифровых каналов, мобильного взаимодействия, маркетинговых коммуникаций и работы с клиентскими данными~\cite{chernova2022, ziolkowska2021}.

Для кафе цифровизация охватывает кассовые и учетные функции, клиентские точки контакта, сервисные сценарии и аналитику. При этом тематическое кафе решает более сложную задачу, поскольку цифровой слой должен поддерживать атмосферу заведения, события, повторные коммуникации, персонализацию и накопление первичной клиентской базы. В такой постановке особое значение приобретает клиентский цифровой контур, то есть совокупность интерфейсов, сценариев, данных и интеграций, сопровождающих гостя до визита, во время посещения и после него. Подходы сервисного дизайна связывают качество услуги с последовательностью точек контакта и совместным созданием ценности между клиентом и организацией~\cite{stickdorn2018, chathoth2016}.

Научная значимость темы усиливается исследованиями клиентской вовлеченности. В работах Р.~Дж.~Броуди, Ш.~Д.~Вивек, А.~Пансари и их соавторов вовлеченность рассматривается как психологическое состояние и форма участия клиента, включающая когнитивные, эмоциональные, поведенческие и социальные проявления~\cite{brodie2011, vivek2012, pansari2017}. Следовательно, цифровой контур тематического кафе может создавать стимулы и точки контакта, которые проявляются в бронированиях, отзывах, повторных обращениях, рекомендациях и данных для дальнейшей CRM-логики.

Для тематического заведения важна также связь цифрового взаимодействия с переживаемым опытом. В терминах экономики впечатлений ценность формируется через продукт, событие, среду и символику бренда~\cite{pine1998}, а геймификация может использоваться как инструмент поддержки роли гостя, прогресса, мотивации и связи цифрового сценария с реальным посещением~\cite{deterding2011, xu2017}. Одновременно проектирование такого контура ограничено отсутствием исторической клиентской базы, ресурсами сопровождения и требованиями законодательства о персональных данных~\cite{fz149, fz152}. Поэтому актуальность настоящей работы определяется необходимостью научно обоснованного проектирования клиентского цифрового контура нового тематического кафе.

\textbf{Объектом исследования} выступают процессы автоматизации и цифровой трансформации деятельности предприятия общественного питания.

\textbf{Предметом исследования} являются методы, модели и инструментальные средства проектирования и разработки клиентского цифрового контура тематического кафе, предназначенного для поддержки клиентского пути, сбора и обработки данных, бронирования, обратной связи и тематических сценариев вовлечения.

\textbf{Хронологические рамки.} Исследование охватывает 2021--2026~гг.; в этот период рассматриваются современные подходы к цифровизации общественного питания, клиентскому взаимодействию, геймификации и обработке персональных данных.

\textbf{Географические рамки.} Исследование ограничено российским рынком общественного питания и правовым пространством Российской Федерации. Зарубежные решения и исследования используются как сравнительный материал для выявления возможностей клиентского цифрового контура.

\textbf{Цель выпускной квалификационной работы} состоит в разработке клиентского цифрового контура для нового тематического кафе (на примере ООО <<\_\_\_\_>>), обеспечивающего поддержку клиентского пути, учет бронирований, формирование первичной клиентской базы, обработку обратной связи и реализацию тематических сценариев вовлечения.

Для реализации поставленной цели в работе были сформулированы и решены следующие \textbf{задачи}:
\begin{enumerate}
    \item Провести анализ теоретических основ цифровизации предприятий общественного питания, рыночных решений, клиентского пути, психологических оснований вовлеченности и нормативно-правовой базы обработки персональных данных.
    \item Выполнить системный анализ предметной области: описать бизнес-процессы в состояниях AS-IS и TO-BE, определить сценарии клиентского цифрового контура и сформировать спецификацию функциональных и нефункциональных требований.
    \item Обосновать технологический способ реализации и стек разработки, спроектировать архитектуру клиент-серверного приложения и логическую структуру реляционной базы данных.
    \item Программно реализовать базовые подсистемы клиентского цифрового контура: интерактивное меню, модуль онлайн-бронирования мероприятий, сбор обратной связи и автоматизированное рабочее место (АРМ) администратора.
    \item Спроектировать и интегрировать тематические сценарии вовлечения, включая элементы геймификации и генеративные языковые модели для поддержки повествовательного опыта заведения.
    \item Провести тестирование разработанного программного обеспечения и выполнить оценку экономической эффективности внедрения.
\end{enumerate}

\textbf{Методы исследования.} В работе применялись анализ, синтез, сравнение, обобщение, системный подход, анализ клиентского пути, сравнительный анализ классов программных решений, структурно-функциональное и объектно-ориентированное моделирование, методы инженерии требований, проектирование реляционной модели данных, метод аналогов, расчет операционной экономии, функциональное и кроссбраузерное тестирование.

\textbf{Методологическая и теоретическая база исследования} включает системный анализ, инженерию программного обеспечения, проектирование баз данных, сервисный дизайн, исследования клиентской вовлеченности и пользовательского опыта~\cite{utkin_baldin_2003, radnaeva2019, bobrovskiy2024, stickdorn2018, brodie2011, vivek2012, pansari2017, pine1998, xu2017}.

\textbf{Степень научной разработанности проблемы.} Автоматизированные информационные системы, цифровизация предприятий, сервисный дизайн, клиентская вовлеченность, геймификация и правовое регулирование персональных данных достаточно подробно представлены в научной и учебно-методической литературе~\cite{utkin_baldin_2003, chernova2022, brodie2011, vivek2012, pansari2017, deterding2011, xu2017, kazakevich2021}. Вместе с тем вопросы проектирования клиентского цифрового контура нового тематического кафе, связывающего клиентский путь, данные, психологические механизмы вовлечения, CRM-логику, геймификацию и правовые ограничения в единую проектную модель, остаются недостаточно разработанными.

\textbf{Научная новизна работы} определяется тем, что предложена модель клиентского цифрового контура нового тематического кафе, объединяющая многослойное представление цифровизации предприятия, клиентский путь, психологическую модель цифрового вовлечения, требования к данным, сценарии геймификации и правовые ограничения обработки персональных данных. На основе данной модели разрабатывается программный прототип, связывающий сервисные функции, первичное накопление клиентской базы и тематические сценарии взаимодействия.

\textbf{Практическая значимость работы} заключается в том, что результаты исследования могут быть использованы при создании прототипа информационной системы для ООО <<\_\_\_\_>>. Система предназначена для цифрового представления меню, учета бронирований, сбора обратной связи, администрирования клиентских сценариев, тематического вовлечения гостей и формирования основы для дальнейшего развития CRM-логики нового кафе.

Структурно выпускная квалификационная работа состоит из введения, трёх основных глав, заключения, списка использованных источников и приложений.

Во введении обосновывается актуальность темы, определяются объект, предмет, цель научно-практического исследования, а также формулируются рабочие задачи.

В первой главе раскрыта логика цифровизации кафе. Выделены основные слои цифровизации, проанализированы классы рыночных решений, обоснован клиентский цифровой контур, раскрыта психологическая модель цифрового вовлечения гостя, определены требования тематического опыта и правовые ограничения работы с персональными данными.

Во второй главе приведена организационно-экономическая специфика исследуемого микропредприятия, осуществлено функциональное моделирование бизнес-процессов (в нотациях AS-IS и TO-BE), сформированы системные требования, конкретизированы сценарии клиентского цифрового контура и обоснован выбор технологического способа реализации.

В третьей главе описан процесс проектирования программной архитектуры и базы данных на инфологическом и даталогическом уровнях, проиллюстрирована практическая реализация основных подсистем клиентского цифрового контура и тематических сценариев вовлечения, представлены итоги тестирования и произведен расчет показателей экономической результативности проекта.

В заключении обобщены результаты проведённой работы и сформулированы итоговые выводы о достижении поставленной цели проектирования.
